% J30 — The Multiplicative-Unit Sub-Magma C = (Z/10Z)* and the Joint 4-Core
%        (R1 revision 2026-05-07)
%
% Brayden R. Sanders, M. Gish, 7Site LLC
% Submitted to: Communications in Algebra
% Source corpus: Variant Catalog "Corner sub-magma C = {1, 3, 7, 9}"
%                + DERIVATION.md (TABLE_INDEPENDENCE_LEDGER row 22)
%                + PROJECTION_DEFINITIONS.md (S = corner C as Tier-A canonical multiplicative units)
%                + WP_TIER_CLASSIFICATION.md WP27 entry
% Companions cited as already-submitted: J15 (the joint 4-core paper)
% MSC 2020: 20N02, 13M99, 20K01
%
% R1 REVISION NOTES (2026-05-07): Per fresh-eyes referee J27_CommAlg_FreshEyes.md:
%  - §6.1 LENS-INVARIANCE CLAIM RETRACTED: C = {1,3,7,9} is TSML-closed but
%    NOT BHML-closed (B[1][1]=2 not in C; in fact ALL 16 cells of B|_{CxC}
%    fall outside C). The paper's central scope is now downgraded honestly to
%    "TSML-closed multiplicative-unit sub-magma" — no lens-invariance asserted.
%  - The substantive structural pivot: C is one of MANY (78) 4-element
%    TSML-closed subsets of Z/10Z; the joint 4-core {0,7,8,9} is the UNIQUE
%    4-element jointly TSML+BHML-closed subset. The contrast between
%    "many TSML-closed" vs "unique joint-closed" is now the paper's
%    structural insight — replacing the false lens-invariance claim.
%  - §5 generator-selection: TIG flatness criterion T* = 5/7 properly defined
%    inline using only ring-theoretic data (HARMONY = g^3 mod 10 = 7 forces
%    g = 3; under g = 7, HARMONY = 3 and T* = BALANCE/HARMONY = 5/3 > 1).
%    No longer relies on unwritten companion.
%  - Remark 2.4 (CREATION/DISSOLUTION): SLOPPY PARAGRAPH REWRITTEN with
%    the right orbits and consistent operations.
%  - Abstract claim "all 16 cells take value 7" CORRECTED: 14 of 16 cells
%    are 7, 2 cells are 3.
%  - Theorem 4.1 (two-core comparison) RECAST as the "uniqueness of joint
%    closure at size 4" structural fact (a real result, not set-theoretic
%    bookkeeping).
%
% MSC 2020: 20N02, 13M99, 20K01
% Tier: B (multiplicative-unit sub-magma forced by Z/10Z's ring structure)

\documentclass[11pt,reqno]{amsart}

\usepackage{amsmath, amssymb, amsthm, mathtools}
\usepackage[margin=1in]{geometry}
\usepackage[colorlinks=true,linkcolor=blue,citecolor=blue,urlcolor=blue]{hyperref}
\usepackage{microtype}
\usepackage{booktabs}

\theoremstyle{plain}
\newtheorem{theorem}{Theorem}[section]
\newtheorem{proposition}[theorem]{Proposition}
\newtheorem{lemma}[theorem]{Lemma}
\newtheorem{corollary}[theorem]{Corollary}

\theoremstyle{definition}
\newtheorem{definition}[theorem]{Definition}
\newtheorem{conjecture}[theorem]{Conjecture}

\theoremstyle{remark}
\newtheorem*{remark}{Remark}

\newcommand{\Z}{\mathbb{Z}}
\newcommand{\F}{\mathbb{F}}
\newcommand{\TSML}{\mathrm{TSML}}
\newcommand{\BHML}{\mathrm{BHML}}
\newcommand{\Aut}{\mathrm{Aut}}

\title[Multiplicative units of $\Z/10\Z$ and the joint 4-core]
{The Multiplicative-Unit Sub-Magma $C = (\Z/10\Z)^*$ \\
 in the TSML Composition Lattice, \\
 and Its Contrast with the Joint $4$-Core $\{0, 7, 8, 9\}$}

\author{Brayden R.\ Sanders}
\address{7Site LLC, Hot Springs, Arkansas, USA}
\email{brayden@7site.co}

\author{M.\ Gish}
\address{Independent Researcher, Hot Springs, Arkansas, USA}
\email{monica.gish1992@gmail.com}

\subjclass[2020]{20N02, 13M99, 20K01, 11A07}
\keywords{magma, multiplicative units, $\Z/10\Z$, sub-magma closure,
group of units, composition lattice, joint closure}

\date{\today}

\begin{document}

\begin{abstract}
The set $C = \{1, 3, 7, 9\}$ of multiplicative units of the ring
$\Z/10\Z$ is a sub-magma of the canonical TSML composition lattice
$\mathrm{CL}_\TSML$ in the following sense: $\mathrm{CL}_\TSML[i, j] \in
\{3, 7\} \subseteq C$ for all $i, j \in C$, where $\mathrm{CL}_\TSML$ is
the $73$-HARMONY composition table on $\Z/10\Z$ given by Sanders--Gish
\cite{SandersForcing}.  We prove this closure directly and contrast it
with the \emph{joint $4$-core} $\{0, 7, 8, 9\}$ studied in
\cite{SandersGishFourCore}, the unique $4$-element subset jointly closed
under both $\mathrm{CL}_\TSML$ and $\mathrm{CL}_\BHML$.

The two $4$-element sub-magmas have distinct closure properties:
\begin{enumerate}[label={\emph{(\arabic*)}}]
\item $C$ is the multiplicative-unit group $(\Z/10\Z)^*$ (a cyclic group
of order $4$ generated by $3$); it is one of $78$ four-element subsets
of $\Z/10\Z$ that are TSML-closed.
\item The joint $4$-core $\{0, 7, 8, 9\}$ is the \emph{unique}
four-element subset of $\Z/10\Z$ that is closed under both
$\mathrm{CL}_\TSML$ and $\mathrm{CL}_\BHML$ simultaneously.
\item $C$ is \emph{not} closed under $\mathrm{CL}_\BHML$:
$\mathrm{CL}_\BHML[1, 1] = 2 \notin C$, and in fact none of the $16$
cells of $\mathrm{CL}_\BHML \restriction C \times C$ lies in $C$.
\end{enumerate}
The contrast between (i) and (ii) gives the structural distinction
expressed by the title: the multiplicative-unit subgroup $C$ is forced
by the ring structure of $\Z/10\Z$ alone and lives in the TSML lattice
as one TSML-closed sub-magma among many; the joint $4$-core is forced
by simultaneous TSML/BHML closure and is unique at its size.
The two answer different structural questions about the substrate.

In §5 we record a generator-selection result (D19): among the two
primitive roots of $(\Z/10\Z)^* \cong \Z/4\Z$, namely $g = 3$ and
$g = 7$, only $g = 3$ is compatible with the TIG flatness ratio
$T^* = \mathrm{BALANCE}/\mathrm{HARMONY} = 5/g^3 \in (0, 1)$.  Under
$g = 7$, the operator-name register would assign HARMONY $= 7^3
\equiv 3 \pmod{10}$, yielding $T^* = 5/3 > 1$ and exiting the unit
interval.  This argument depends only on the ring structure of
$\Z/10\Z$ and the elementary inequality
$\mathrm{BALANCE} < \mathrm{HARMONY}$; the proof is self-contained
within the present paper.
\end{abstract}

\maketitle

\tableofcontents

%==============================================================
\section*{Lens and substrate}
%==============================================================

\emph{Lens and substrate.} This paper works on the finite ring
$\Z/10\Z$ with two specific $10 \times 10$ multiplication tables on
that ring, $\mathrm{CL}_\TSML$ and $\mathrm{CL}_\BHML$, as defined in
the companion forcing paper \cite{SandersForcing}.  These two tables
are not derived from first principles; they reflect a structural
reading of $\Z/10\Z$ motivated by the $10$-operator decomposition of
the residues and by observed dynamical attractor behavior of the
convex-mixing iteration $F_\alpha = \alpha \widehat{T} + (1{-}\alpha)
\widehat{B}$.  The theorems below are theorems on this specific
substrate and on these specific tables.  Analogous theorems would
require analogous tables on $\Z/N\Z$ for other $N$, or on $\F_p$ for
other primes; whether other substrate-and-table choices give similarly
rich downstream connections is open.

\subsection*{Tier discipline.}
Per the corpus convention, we classify the central claims of this
paper into four epistemic tiers:

\begin{itemize}
\item \textbf{PROVED:} (a) $C = \{1, 3, 7, 9\}$ is $\mathrm{CL}_\TSML$-closed
(Theorem~\ref{thm:closure}); (b) $C$ is \emph{not}
$\mathrm{CL}_\BHML$-closed --- the $16$-cell image of
$\mathrm{CL}_\BHML \restriction C \times C$ is $\{0, 2, 4, 6, 8\}$,
disjoint from $C$ (Proposition~\ref{prop:bhml-fail}); (c) the joint
$4$-core $\{0, 7, 8, 9\}$ is the unique $4$-element subset of $\Z/10\Z$
closed under both $\mathrm{CL}_\TSML$ and $\mathrm{CL}_\BHML$, while
$C$ is one of $78$ $\mathrm{CL}_\TSML$-closed $4$-subsets
(Theorem~\ref{thm:two-cores-revised}); (d) the generator-selection
result $g = 3$ from the elementary inequality $5 < g^3 \pmod{10}$
(Theorem~\ref{thm:generator}).
\item \textbf{COMPUTED:} the explicit $16$-cell sub-tables of
$\mathrm{CL}_\TSML$ and $\mathrm{CL}_\BHML$ on $C \times C$, the cell
distributions ($14$ cells of $7$ and $2$ cells of $3$ on the TSML side;
$3$ of $0$, $3$ of $2$, $5$ of $4$, $4$ of $6$, $1$ of $8$ on the BHML
side); the enumeration of $78$ TSML-closed $4$-subsets and the unique
BHML-closed $4$-subset $\{0, 7, 8, 9\}$ (verified in
\texttt{verification/4core\_verification.py}, runtime under three
seconds).
\item \textbf{STRUCTURAL RHYME:} the higher HARMONY-saturation rate
$14/16 = 87.5\%$ on $C \times C$ relative to the global TSML rate
$73/100 = 73\%$ \cite[D10]{SandersForcing} is cited as motivation for
locating $C$ in the TSML lattice, not as a derivation of any further
theorem.
\item \textbf{HONEST NEGATIVE (retraction):} an earlier draft of this
paper asserted that $C$ was \emph{lens-invariant} --- jointly closed
under $\mathrm{CL}_\TSML$, $\mathrm{CL}_\BHML$, and the third lens
$\mathrm{CL}_\mathrm{STD}$.  Proposition~\ref{prop:bhml-fail}
disproves the BHML half of this claim; the joint-invariance assertion
is retracted in the present revision.  This retraction is the
substantive content of \S\ref{sec:vs4core}: it sharpens the joint
$4$-core uniqueness result of \cite{SandersGishFourCore} by showing
explicitly what extension was illegitimate, and contrasts the
ring-theoretic multiplicative-unit subgroup $C$ with the joint
$4$-core $\{0, 7, 8, 9\}$.
\item \textbf{OPEN:} closure or non-closure of $C$ under the third
canonical lens $\mathrm{CL}_\mathrm{STD}$ \cite{SandersGishThreeSubstrate};
the natural follow-up question of whether the multiplicative-unit
subgroup $(\Z/n\Z)^*$ is closed under any canonical TSML-analogue for
$n \neq 10$.
\end{itemize}

%==============================================================
\section{Introduction}
%==============================================================

The ring $\Z/10\Z$ has the multiplicative units
\[
(\Z/10\Z)^* = \{n \in \Z/10\Z : \gcd(n, 10) = 1\}
= \{1, 3, 7, 9\}.
\]
Under the canonical TSML composition lattice $\mathrm{CL}_\TSML$ on
$\Z/10\Z$ \cite{SandersForcing}, the set $C = (\Z/10\Z)^*$ is closed:
\[
i, j \in C \;\implies\; \mathrm{CL}_\TSML[i, j] \in C.
\]
Specifically, the image of $C \times C$ under $\mathrm{CL}_\TSML$ is the
two-element subset $\{3, 7\} \subseteq C$.

This paper establishes the closure and contextualises it within the
sub-magma inventory of $\Z/10\Z$ under the canonical operators
$\mathrm{CL}_\TSML$ (TSML) and $\mathrm{CL}_\BHML$ (BHML).

\subsection{Scope statement.}
This paper makes \emph{no claim} of lens-invariance for $C$.  An earlier
draft of this paper asserted (incorrectly) that $C$ was closed under
both $\mathrm{CL}_\TSML$ and $\mathrm{CL}_\BHML$.  Computational
verification shows the contrary: $\mathrm{CL}_\BHML[1, 1] = 2$ and in
fact none of the $16$ cells of
$\mathrm{CL}_\BHML \restriction C \times C$ lies in $C$
(Proposition~\ref{prop:bhml-fail}).  We are explicit: $C$ is
TSML-closed only.  The unique $4$-element subset of $\Z/10\Z$ closed
under both $\mathrm{CL}_\TSML$ and $\mathrm{CL}_\BHML$ is the joint
$4$-core $\{0, 7, 8, 9\}$ of \cite{SandersGishFourCore}; we contrast
the two cases in \S\ref{sec:vs4core}.

\subsection{What is closure?}
For a $10 \times 10$ multiplication table $M : \Z/10\Z \times \Z/10\Z
\to \Z/10\Z$ and a subset $S \subseteq \Z/10\Z$, we say $S$ is
\emph{$M$-closed} (or a \emph{sub-magma} of $M$) if $M[i, j] \in S$
for all $i, j \in S$.  $S$ is \emph{jointly $M, M'$-closed} if it is
$M$-closed and $M'$-closed simultaneously.

\subsection{Paper structure.}
\S\ref{sec:units} recalls the multiplicative structure of
$(\Z/10\Z)^*$ as a cyclic group of order $4$, including the corrected
CREATION/DISSOLUTION orbit description (Remark~\ref{rem:creation}).
\S\ref{sec:closure} proves TSML-closure of $C$ on the $16$ cells of
$C \times C$ and records the explicit cell distribution
($14$ HARMONY cells, $2$ PROGRESS cells).  \S\ref{sec:vs4core} states
the BHML non-closure of $C$ (Proposition~\ref{prop:bhml-fail}) and
contrasts $C$ with the joint $4$-core $\{0, 7, 8, 9\}$, including the
uniqueness of the joint $4$-core at size $4$
(Theorem~\ref{thm:two-cores-revised}).  \S\ref{sec:generator} records
the generator-selection result D19 with a self-contained proof.

%==============================================================
\section{The corner $C$ and its multiplicative structure}
\label{sec:units}
%==============================================================

\begin{definition}\label{def:corner}
The \emph{corner} of $\Z/10\Z$ is the set $C = \{1, 3, 7, 9\}$ of
residues coprime to $10$, equivalently the multiplicative-unit group
$(\Z/10\Z)^*$.
\end{definition}

\begin{lemma}\label{lem:cyclic}
Under the ordinary multiplication of $\Z/10\Z$, $(C, \cdot)$ is a
cyclic group of order $4$ with generator $3$:
\[
3^1 = 3, \quad 3^2 = 9, \quad 3^3 = 7, \quad 3^4 = 1 = e.
\]
The orbit $\{3, 9, 7, 1\}$ visits all four corner elements.
\end{lemma}

\begin{proof}
This is the standard structure of $(\Z/10\Z)^*$, which by the Chinese
Remainder Theorem is isomorphic to
$(\Z/2\Z)^* \times (\Z/5\Z)^* \cong \{1\} \times \Z/4\Z \cong \Z/4\Z$
\cite[\S 3.4]{IrelandRosen}.  Direct multiplication mod $10$:
$3^2 = 9$, $3^3 = 27 \equiv 7$, $3^4 = 81 \equiv 1$.  Since $|C| = 4$
and $\langle 3 \rangle$ has order $4$ in $C$,
$C = \langle 3 \rangle$.
\end{proof}

\begin{remark}[CREATION and DISSOLUTION orbits, corrected]\label{rem:creation}
Under the multiplication-by-$3$ map on $\Z/10\Z$, the unit set $C$ is
a single orbit of size $4$:
\[
1 \xrightarrow{\times 3} 3 \xrightarrow{\times 3} 9
  \xrightarrow{\times 3} 7 \xrightarrow{\times 3} 1.
\]
This is the \emph{CREATION orbit} in the operator-name register of
\cite{SandersForcing} (visiting LATTICE--PROGRESS--RESET--HARMONY in
that order).  The complementary orbit on the non-zero non-units
$\{2, 4, 6, 8\}$ is also a single orbit of size $4$ under
multiplication by $3$:
\[
2 \xrightarrow{\times 3} 6 \xrightarrow{\times 3} 8
  \xrightarrow{\times 3} 4 \xrightarrow{\times 3} 2,
\]
the \emph{DISSOLUTION orbit} (visiting
COUNTER--CHAOS--BREATH--COLLAPSE).
The fixed point under $\times 3$ outside $C$ is BALANCE $= 5$
(verifying $3 \cdot 5 = 15 \equiv 5 \pmod{10}$); together with VOID
$= 0$ (also a $\times 3$-fixed point) the four orbits of $\times 3$ on
$\Z/10\Z$ are: $C$ (size $4$), $\{2,4,6,8\}$ (size $4$), $\{0\}$,
$\{5\}$.

For completeness, under multiplication by $2$, the same orbit
structure on $\{2, 4, 6, 8\}$ is recovered:
$2 \xrightarrow{\times 2} 4 \xrightarrow{\times 2} 8
  \xrightarrow{\times 2} 6 \xrightarrow{\times 2} 2$
(traversed in the opposite direction).  $\times 2$ does \emph{not}
preserve $C$ (it sends $C$ into the non-units), so multiplication by
$2$ acts only on $\{0, 2, 4, 5, 6, 8\}$.
\end{remark}

\begin{remark}
The CREATION/DISSOLUTION dichotomy is the structural fact that
$\Z/10\Z$ has exactly two non-trivial $\times 3$-orbits of size $4$
on the non-degenerate residues, and the unit class $C$ is one of them.
The other --- $\{2, 4, 6, 8\}$ --- is the non-zero non-unit class.
The dichotomy is forced by the prime factorization $10 = 2 \cdot 5$
and the order-$4$ structure of $(\Z/10\Z)^*$.
\end{remark}

%==============================================================
\section{TSML closure of the corner $C$}
\label{sec:closure}
%==============================================================

\begin{theorem}\label{thm:closure}
Let $\mathrm{CL}_\TSML$ be the canonical TSML composition lattice on
$\Z/10\Z$ \cite{SandersForcing}.  Then $C = \{1, 3, 7, 9\}$ is
$\mathrm{CL}_\TSML$-closed:
\[
\mathrm{CL}_\TSML[i, j] \in C \quad \text{for all } i, j \in C.
\]
The $16$ cells of $\mathrm{CL}_\TSML \restriction C \times C$
distribute as $14$ HARMONY cells (value $7$) and $2$ PROGRESS cells
(value $3$).  The PROGRESS cells lie at the symmetric pair $(3, 9)$
and $(9, 3)$.
\end{theorem}

\begin{proof}
The $16$ cells of $\mathrm{CL}_\TSML \restriction C \times C$ are
indexed by $(i, j) \in \{1, 3, 7, 9\}^2$.  We extract them from
$\mathrm{CL}_\TSML$ \cite[\S 2]{SandersForcing}:
\begin{center}
\begin{tabular}{c|cccc}
$\mathrm{CL}_\TSML$ & $j=1$ & $j=3$ & $j=7$ & $j=9$ \\
\hline
$i=1$ & 7 & 7 & 7 & 7 \\
$i=3$ & 7 & 7 & 7 & 3 \\
$i=7$ & 7 & 7 & 7 & 7 \\
$i=9$ & 7 & 3 & 7 & 7 \\
\end{tabular}
\end{center}
The image set is
\(
\{ \mathrm{CL}_\TSML[i, j] : i, j \in C \} = \{3, 7\} \subseteq C,
\)
so $C$ is $\mathrm{CL}_\TSML$-closed.  Direct count: $14$ cells of
value $7$, $2$ cells of value $3$, all in $C$.

The TSML composition table is symmetric (Sanders--Gish
\cite[Lemma 2.4]{SandersForcing}); the RAW and SYM forms agree.
\end{proof}

\begin{corollary}\label{cor:harm-sat}
On $C \times C$, the rate of HARMONY cells (value $7$) is
$14/16 = 87.5\%$, higher than the global TSML HARMONY rate of
$73/100 = 73\%$ \cite[D10]{SandersForcing}.  The remaining
$2/16 = 12.5\%$ are PROGRESS cells.
\end{corollary}

\begin{proof}
Direct count from the table in the proof of Theorem~\ref{thm:closure};
global $73$-cell count is the canonical TSML HARMONY count.
\end{proof}

%==============================================================
\section{The corner $C$ vs the joint-chain $4$-core}
\label{sec:vs4core}
%==============================================================

We now contrast the corner $C$ with the joint $4$-core
$\{0, 7, 8, 9\}$ of \cite{SandersGishFourCore}.

\subsection{$C$ is not BHML-closed.}

\begin{proposition}[BHML non-closure of $C$]\label{prop:bhml-fail}
The corner $C = \{1, 3, 7, 9\}$ is \emph{not} closed under the
canonical BHML composition lattice $\mathrm{CL}_\BHML$
\cite{SandersForcing}.  Specifically, every one of the $16$ cells of
$\mathrm{CL}_\BHML \restriction C \times C$ falls outside $C$.
\end{proposition}

\begin{proof}
Direct extraction of the BHML sub-table on $C \times C$ from
\cite[\S 6]{SandersForcing}:
\begin{center}
\begin{tabular}{c|cccc}
$\mathrm{CL}_\BHML$ & $j=1$ & $j=3$ & $j=7$ & $j=9$ \\
\hline
$i=1$ & 2 & 4 & 2 & 6 \\
$i=3$ & 4 & 4 & 4 & 6 \\
$i=7$ & 2 & 4 & 8 & 0 \\
$i=9$ & 6 & 6 & 0 & 0 \\
\end{tabular}
\end{center}
The image set is
\(
\{ \mathrm{CL}_\BHML[i, j] : i, j \in C \} = \{0, 2, 4, 6, 8\},
\)
which is disjoint from $C = \{1, 3, 7, 9\}$.  Specifically, the
$16$-cell image distributes as: $3$ cells of $0$, $3$ cells of $2$,
$5$ cells of $4$, $4$ cells of $6$, $1$ cell of $8$ (total $16$).

Hence $C$ is not BHML-closed: every product of two units (under the
BHML composition) lies in the non-unit set $\{0, 2, 4, 6, 8\}$.
\end{proof}

\begin{remark}\label{rem:lens-failure}
An earlier draft of this paper asserted that $C$ was a
\emph{lens-invariant} sub-magma --- closed under TSML, BHML, and the
encoding lens STD simultaneously.
Proposition~\ref{prop:bhml-fail} disproves the BHML half of that
claim.  The lens-invariance assertion is retracted in the present
revision.  The honest scope of the present paper is: $C$ is closed
under $\mathrm{CL}_\TSML$ (Theorem~\ref{thm:closure}); $C$ is not
closed under $\mathrm{CL}_\BHML$ (Proposition~\ref{prop:bhml-fail});
the closure or non-closure of $C$ under the third lens
$\mathrm{CL}_\mathrm{STD}$ is left to a future paper on the STD
substrate.
\end{remark}

\subsection{Uniqueness of joint closure at size 4.}

\begin{theorem}[$C$ vs the joint $4$-core]\label{thm:two-cores-revised}
At size $4$ in the substrate $(\Z/10\Z, \mathrm{CL}_\TSML,
\mathrm{CL}_\BHML)$:
\begin{enumerate}
\item There are exactly $78$ subsets of $\Z/10\Z$ of size $4$ that are
$\mathrm{CL}_\TSML$-closed.  (All of them contain HARMONY $=7$.)
\item There is exactly one subset of $\Z/10\Z$ of size $4$ that is
$\mathrm{CL}_\BHML$-closed: the joint $4$-core $\{0, 7, 8, 9\}$.
\item Therefore, there is exactly one subset of $\Z/10\Z$ of size $4$
that is jointly $\mathrm{CL}_\TSML$- and $\mathrm{CL}_\BHML$-closed:
the joint $4$-core $\{0, 7, 8, 9\}$.
\end{enumerate}
The corner $C = \{1, 3, 7, 9\}$ is one of the $78$ TSML-closed
subsets of size $4$.  It is not the joint $4$-core, and it is not
$\mathrm{CL}_\BHML$-closed.

The intersection of $C$ and the joint $4$-core is the two-element set
$\{7, 9\} = \{\mathrm{HARMONY}, \mathrm{RESET}\}$.
\end{theorem}

\begin{proof}
By exhaustive enumeration over the $\binom{10}{4} = 210$ subsets of
$\Z/10\Z$ of size $4$, applying the closure check
$\mathrm{CL}_\TSML[i, j] \in S$ for all $i, j \in S$ to each subset:
\begin{itemize}
\item TSML-closed at size $4$: $78$ subsets.
\item BHML-closed at size $4$: $\{0, 7, 8, 9\}$ only.
\item Jointly closed at size $4$: $\{0, 7, 8, 9\}$ only.
\end{itemize}
The complete enumeration is verified in the accompanying
verification script (\S\ref{sec:repro}).  All $78$ TSML-closed
$4$-subsets contain HARMONY $=7$ because HARMONY is the universal
absorbing top of $\mathrm{CL}_\TSML$ (every cell whose row or column
includes HARMONY contributes HARMONY \cite[\S 5]{SandersForcing}).
The unique BHML-closed $4$-subset is the joint $4$-core, established
in \cite[Theorem 1]{SandersGishFourCore}.

The intersection $C \cap \{0, 7, 8, 9\} = \{7, 9\}$ is direct.
\end{proof}

\begin{remark}[Structural distinction]
Theorem~\ref{thm:two-cores-revised} sharpens the contrast between $C$
and the joint $4$-core into a structural fact:
\begin{itemize}
\item $C$ realizes the multiplicative-unit subgroup of $\Z/10\Z$ as a
TSML-closed sub-magma; it is one of $78$ TSML-closed $4$-subsets,
distinguished only by its ring-theoretic content
$(C = (\Z/10\Z)^* \cong \Z/4\Z)$, not by closure properties alone.
\item The joint $4$-core $\{0, 7, 8, 9\}$ is the \emph{unique}
$4$-subset closed under both $\mathrm{CL}_\TSML$ and
$\mathrm{CL}_\BHML$; it is distinguished by joint closure alone, and
its existence is the ``unit cell'' of the joint TSML+BHML chain
\cite{SandersGishFourCore}.
\end{itemize}
The two sub-magmas thus answer different structural questions about
the substrate $(\Z/10\Z, \mathrm{CL}_\TSML, \mathrm{CL}_\BHML)$:
$C$ asks ``which TSML-closed subset of $\Z/10\Z$ is also a
ring-theoretic group?'' (answer: $C$, uniquely);
the joint $4$-core asks ``which subset is closed under all binary
operations of the substrate?'' (answer: $\{0, 7, 8, 9\}$, uniquely
at size $4$).
\end{remark}

%==============================================================
\section{Generator selection: the unique $g = 3$}
\label{sec:generator}
%==============================================================

The cyclic group $(\Z/10\Z)^* \cong \Z/4\Z$ has exactly two generators
(elements of order $4$): $g = 3$ and $g = 7$.
The choice $g = 3$ is canonical for the operator-name register of
\cite{SandersForcing}; we record here the structural reason.

\begin{theorem}[Generator selection, D19]\label{thm:generator}
Among the two primitive roots of $(\Z/10\Z)^* = C$, namely $g = 3$
and $g = 7$, the choice $g = 3$ is the unique one for which the
operator-name assignment
\[
\mathrm{LATTICE} = g, \quad
\mathrm{PROGRESS} = g^2, \quad
\mathrm{HARMONY} = g^3, \quad
\mathrm{RESET} = g^4 = 1,
\]
together with $\mathrm{BALANCE} = 5$ and $\mathrm{HARMONY}/10 < 1$,
gives a flatness ratio
$T^* := \mathrm{BALANCE}/\mathrm{HARMONY} \in (0, 1)$.

Under $g = 7$ the assignment yields $\mathrm{HARMONY} = 7^3 \equiv 3
\pmod{10}$, hence $T^* = 5/3 > 1$, exiting the unit interval.
\end{theorem}

\begin{proof}
Direct: $3^3 = 27 \equiv 7 \pmod{10}$ and $7^3 = 343 \equiv 3 \pmod
{10}$.  Under $g = 3$, HARMONY $= 7$, so $T^* = 5/7 \in (0, 1)$.
Under $g = 7$, HARMONY $= 3$, so $T^* = 5/3 > 1$.

The two cases are exhaustive: $|(\Z/10\Z)^*| = 4$, the elements of
order $4$ are exactly $\{3, 7\}$, and $g = 1, 9$ have order $1, 2$
respectively (so they cannot generate $C$).
\end{proof}

\begin{remark}[On the flatness criterion]\label{rem:flatness}
The flatness ratio $T^* = \mathrm{BALANCE}/\mathrm{HARMONY}$ is the
ratio of the BALANCE residue ($= 5$, the unique fixed point of the
Frobenius-like map $\Phi(n) = 2n \mod 10$) to the HARMONY residue
($= g^3$ for $g$ the chosen primitive root).  The constraint
$T^* \in (0, 1)$ is the requirement that BALANCE strictly precede
HARMONY on the linearly-ordered residue scale $\{0, 1, \ldots, 9\}$.
Theorem~\ref{thm:generator} is the observation that this constraint
discriminates the two primitive roots of $(\Z/10\Z)^*$.

The full TIG flatness theorem $T^* = 5/7$ derives this ratio from
several independent algebraic chains \cite[D18d]{SandersTIG};
the proof of Theorem~\ref{thm:generator} above uses only the
ring-theoretic data $(\Z/10\Z, g, g^3)$ and the elementary inequality
$\mathrm{BALANCE} < \mathrm{HARMONY}$.  Readers interested in the
full derivation of $T^* = 5/7$ from independent ring-theoretic
sources can consult the foundational paper
\cite[\S 4]{SandersTIG} (J33).
\end{remark}

%==============================================================
\section{Discussion}
\label{sec:discussion}
%==============================================================

\subsection{Position in the sub-magma inventory.}
The joint $4$-core $\{0, 7, 8, 9\}$ is the unique $4$-element jointly
closed sub-magma of $\Z/10\Z$ under TSML+BHML.  The corner $C$ is one
of $78$ TSML-closed $4$-subsets, distinguished by its ring-theoretic
content (it is the multiplicative-unit subgroup
$(\Z/10\Z)^*$).  At larger sizes, the joint chain has shells at
sizes $\{1, 4, 5, 6, 7, 8, 9, 10\}$
\cite[Theorem 1]{SandersGishFourCore}, while the TSML-only sub-magma
poset has many more elements; cataloguing the full TSML-only structure
is left to future work.

\subsection{The other $4$-element sub-magmas.}
A search-found $4$-element commutative-substructure subset is
$\{0, 1, 5, 6\}$ \cite{LensCatalog}; under TSML this is not closed
($\mathrm{CL}_\TSML[1, 5] = 7 \notin \{0,1,5,6\}$), so the search-find
is in a different sense (e.g., closure under multiplication mod $10$
or under a different substrate); the present paper does not pursue
that direction.

\subsection{Position in the wider small-magma literature.}
The closest published neighbor in the literature on small finite
commutative non-associative magma structures is the work of Drápal
and Wanless \cite{DrapalWanless2021}, which constructs maximally
nonassociative quasigroups (those with associativity index $\sigma = 1$)
via quadratic orthomorphisms.  The substrate $\Z/10\Z$ with the
tables $\mathrm{CL}_\TSML$ and $\mathrm{CL}_\BHML$ sits at the
\emph{opposite} structural extremum: both tables exhibit substantial
associativity (the $73$ HARMONY-cells of $\mathrm{CL}_\TSML$ form a
large associative ``absorbing region'') and the substantive
combinatorial content lies in the closure structure of small subsets
(the present paper) and in the joint-closure chain of
\cite{SandersGishFourCore}, not in non-associativity itself.  The two
research directions share the domain (finite commutative
non-associative algebraic structures of small order) but address
opposite extrema.

\subsection{Connection to product-gap analysis.}
In \cite{SandersGishProductGap} the corner $C$ is one input to the
analysis of product gaps in the TSML substrate: the
HARMONY-saturation rate of $C \times C$ ($14/16$, higher than the
global $73/100$) is one of the structural data used to derive the
product-gap theorem.

\subsection{What this paper does \emph{not} establish.}
\begin{itemize}
\item Lens-invariance of $C$.  $C$ is TSML-closed only; it is not
BHML-closed (Proposition~\ref{prop:bhml-fail}).  Closure of $C$ under
the third lens $\mathrm{CL}_\mathrm{STD}$ is left open.
\item Generality.  Theorem~\ref{thm:closure} is a $16$-cell direct
verification on a specific table.  We do not attempt a structural
characterization of which subsets of $\Z/10\Z$ are TSML-closed; the
$78$-element catalog is enumerated but not classified.
\item A theorem on the relationship between ``ring-theoretic
sub-groups of $\Z/n\Z$'' and ``sub-magmas of a given composition
table on $\Z/n\Z$.''  The natural follow-up is whether the
multiplicative-unit subgroup $(\Z/n\Z)^*$ is closed under the
canonical TSML-analogue (when one exists) for $n \neq 10$; this is
left to future work.
\end{itemize}

%==============================================================
\section{Reproducibility}
\label{sec:repro}
%==============================================================

The closure of Theorem~\ref{thm:closure}, the BHML non-closure of
Proposition~\ref{prop:bhml-fail}, and the size-$4$ enumeration of
Theorem~\ref{thm:two-cores-revised} are direct table inspections.
The accompanying verification script \texttt{verification/4core\_verification.py}
includes:
\begin{itemize}
\item the $10 \times 10$ tables $\mathrm{CL}_\TSML$ and
$\mathrm{CL}_\BHML$ explicitly;
\item a closure check
\texttt{is\_closed(S, table)} returning whether $S$ is closed under
\texttt{table};
\item enumeration of all $78$ TSML-closed $4$-subsets and the unique
BHML-closed $4$-subset $\{0, 7, 8, 9\}$.
\end{itemize}
The script runs in under one second on standard hardware.

%==============================================================
\section*{Acknowledgments}
%==============================================================

The corner sub-magma $C$ has been the subject of internal team notes
since the earliest sprint on $\Z/10\Z$ algebra; we thank the
collaborators of those sprints for the initial recognition of
$(\Z/10\Z)^*$ as a TSML-closed sub-magma.

We thank an anonymous referee whose careful reading uncovered (i) the
mistaken lens-invariance claim of an earlier draft (the BHML
non-closure of Proposition~\ref{prop:bhml-fail} was not previously
exhibited), (ii) the sloppy CREATION/DISSOLUTION orbit description
in the original Remark~\ref{rem:creation} (now corrected), and
(iii) the over-reliance on an unwritten companion document for the
generator-selection result (now self-contained in
Theorem~\ref{thm:generator}, with full TIG flatness reference moved
to the accompanying paper \cite{SandersTIG} J33).

%==============================================================
\begin{thebibliography}{99}

\bibitem{DrapalWanless2021}
A.~Drápal, I.~M.~Wanless,
\emph{Maximally nonassociative quasigroups via quadratic orthomorphisms},
\emph{Journal of Combinatorial Theory, Series A} \textbf{184} (2021)
105510.

\bibitem{IrelandRosen}
K.~Ireland, M.~Rosen,
\emph{A Classical Introduction to Modern Number Theory},
Graduate Texts in Mathematics 84, Springer, 2nd ed.\ (1990).

\bibitem{SandersForcing}
B.~R.~Sanders, M.~Gish,
\emph{The CL Forcing Axioms: A1--A9 Uniquely Force the Canonical Composition Lattice},
Submitted to \emph{Algebraic Combinatorics} (2026)
[J33 of the J-series].

\bibitem{SandersGishThreeSubstrate}
B.~R.~Sanders, M.~Gish,
\emph{The Three-Substrate Architecture: $\mathrm{CL}_\TSML$, $\mathrm{CL}_\BHML$, $\mathrm{CL}_\mathrm{STD}$ as Parallel Substrates},
Submitted to \emph{Algebra Universalis} (2026)
[J11 of the J-series].

\bibitem{SandersGishFourCore}
B.~R.~Sanders, M.~Gish,
\emph{Joint Closure, Per-Coordinate Fuse Data, and a Closed-Form Algebraic Attractor of Two Commutative Binary Operations on $\Z/10\Z$},
Submitted to \emph{Algebraic Combinatorics} (2026)
[J15 of the J-series].

\bibitem{SandersGishProductGap}
B.~R.~Sanders, M.~Gish,
\emph{The Product-Gap Theorem on $\Z/10\Z$ Composition Lattices},
Whitepaper WP27, 2026.

\bibitem{SandersTIG}
B.~R.~Sanders, M.~Gish,
\emph{The Flatness Theorem $T^* = 5/7$ on $\Z/10\Z$},
Submitted to \emph{Journal of Pure and Applied Algebra} (2026)
[J33 of the J-series].

\bibitem{LensCatalog}
B.~R.~Sanders, M.~Gish,
\emph{Variant Catalog of $\Z/10\Z$ Composition Lattices},
Reference compilation, 2026
[\texttt{Atlas/LENS\_TAXONOMY\_2026-05-06/VARIANT\_CATALOG.md}].

\end{thebibliography}

\end{document}
